% !PRINTED 52
\weberpage{52}

%<<CONT
oder auch, indem man $u$ in $u+2iK'$ verwandelt und die Bezeichnung $\varkappa$ wieder wegläßt:
\weq{13}{\mathrm{sn}\,(u+2iK') = \mathrm{sn}\,u.}
Es hat also die Funktion sn $u$ eine \gesperrt{doppelte Periodizität}. Die eine dieser Perioden, $4K$, ist reell die andere, $2iK$, rein imaginär % [sic] Printed without the mathematically expected prime on K.
 (wenigstens wenn der Modul $\varkappa$ ein positiver echter Bruch ist). Diese doppelte Periodizität ist eine Eigenschaft, die in dem Gebiete der elementaren Funktionen nirgends vorkommt, und so am deutlichsten zeigte, daß man es hier mit einer neuen Funktionenart zu tun hat. Die aus der doppelten Periodizität abgeleiteten Folgerungen sind es zugleich, die, wie die Erfahrung gezeigt hat, weitaus am schnellsten mit befriedigender Strenge in das Innere der Theorie einführen, so daß hier der zweckmäßigste Ausgangspunkt für eine methodische Entwickelung zu suchen ist. Unsere nächsten Betrachtungen werden daher den doppelt periodischen Funktionen im allgemeinen gewidmet sein.
\begin{center}
\rule{3cm}{0.4pt}
\end{center}
